\section{Normative Language and Goals}

The keywords \MUST{}, \MUSTNOT{}, \SHOULD{}, \SHOULDNOT{}, and \MAY{} are normative.

This specification distinguishes between:

\begin{lstlisting}
contract-level convention:
    a pattern written in MeTTa source code

runtime enforcement:
    a check that cannot be bypassed by malicious MeTTa code
\end{lstlisting}

For DeFi safety, all critical protections \MUST{} be runtime-enforced.

\subsection{Goals}

\SCS{} \MUST{} provide:

\begin{enumerate}
  \item one secured state space per deployed contract instance;
  \item public method calls through a contract dispatcher;
  \item deny-by-default method authorization;
  \item functional access control based on method policy;
  \item live capability-registry lookup as the default authority check;
  \item no routine capability-token passing by callers;
  \item guarded grant, revoke, update, and consume operations for registry capabilities;
  \item runtime-derived caller identity;
  \item method-scoped write authority;
  \item storage-level rejection of unauthorized writes;
  \item atomic commit/revert for state, registry, and events;
  \item invariant enforcement before commit;
  \item deterministic execution;
  \item restricted \PeTTa{} evaluator mode for contract code;
  \item auditable authorization decisions and state transitions.
\end{enumerate}

\subsection{Design bias}

The design deliberately prefers simple, live authorization over portable bearer authority:

\begin{lstlisting}
Default:
    method policy + current caller + live registry lookup

Exception:
    signed certificate adapter at protocol edges
\end{lstlisting}

This keeps ordinary contract calls simple and makes revocation, spend tracking, roles, and governance easier to reason about.
